\def\velkostpisma{9}
\def\velkostriadku{12}
\input makra.tex % nacitanie Ivanom pripravenych nastaveni a prikazov
\hyphenation{star-šov-stvo} % rozdelenie slov na konci riadku, treba tu uviest slova, ktore sam nepozna

\spravodaj{9}{2026}


\clanok{{\it OHLIADNUTIE SA ZA LETOM}}
Aj keď sa september už preklápa do svojej druhej polovice, máme pre Vás pripravené ešte správy z~letného obdobia v~našom zbore.


\clanok{Denný detský tábor}
Na začiatku prázdnin sme v~našom cirkevnom zbore zorganizovali denný tábor.
Bola to prvá lastovička po viacerých rokoch. Prípravy bežali už niekoľko mesiacov vopred a priznám sa, že som netušila, ako celé podujatie dopadne. Avšak hneď na začiatku som sa modlila: „Pane, nechávam ti celú réžiu V~Tvojich rukách, lebo ja sa cítim slabá a neschopná.” Pán Boh ma prekvapil hneď na prvom prípravnom stretnutí ešte začiatkom roka. Takmer všetci, ktorých som pozvala do spolupráce, povedali ÁNO. (Ten, kto organizuje akékoľvek podujatie, vie, aké ťažké je získať ochotných dobrovoľníkov…). Nikoho som nepotrebovala presviedčať, prehovárať, jednoducho sa rozhodli ochotne dávať. Medzi vedúcimi na dennom tábore sa tak stretli kazatelia aj mládežníci z~nášho zboru, mladí i starší, temperamentnejšie aj pokojnejšie povahy. Pán Boh nás viedol pri prípravách aj na samotnom tábore. Vytvorili sme jednotný tím, hoci občas sme mali na niektoré veci rozdielne názory. Počas tábora bol každý na svojom mieste a slúžil deťom telom i dušou. Veľmi si vážim ochotu a nadšenie, s~akým sa všetci vedúci rozdávali.

Najlepším zhodnotením tábora boli pre mňa reakcie detí a rodičov. Jedna mamička mi napísala: „Sárka (najmladšia účastníčka tábora) si stále spieva Boh ma miiilujeee, Spasiteľ ma ľúuuubi...“ Pieseň s~takýmto textom sme na tábore nespievali, ale Sárka si z~tábora zobrala to najpodstatnejšie a vytvorila si vlastnú pesničku. Iná mamička mi poslala povzbudzujúcu správu: „...chcela som Vám veľmi pekne
poďakovať za celý tábor, na ktorom sa Riško cítil výborne a hneď sa pýtal, či bude môcť ísť aj na budúci rok. Ďakujeme za Váš čas, za celú organizáciu a
v~neposlednom rade za to, že Riško mohol skrze Váš tábor zas viac spoznať Boha. Vedzte, že Vaša práca má zmysel a sme Vám za ňu veľmi vďační.” Modlím sa, aby to „málo”, čo sme deťom dali, Pán Boh požehnal a rozmnožil.
Som veľmi vďačná tiež za to, že nielen náš zbor, ale mnohé cirkevné zbory na Slovensku organizujú pre deti rôzne podujatia, na ktorých slúžia ochotní bratia a sestry.
\autor{Miriam Kešjarová}


\clanok{Mládežnícko-dorastenecký tábor v~Novej Lehote}
Na začiatku prázdnin sme mali obrovské privilégium stráviť týždeň s~11 dorastencami a 19 mládežníkmi.
Štruktúra táborov je vždy takmer rovnaká: téma, skupinky, hry, plnenie vatier, chvály a táborák. Tento rok však priniesol aj zmeny. Hlavnou zmenou bola výmena kuchárov. Našu vernú dvojicu Hanku~Š. a Zuzku~P. vystriedali Miško Vrábel a Jany Kerekréty. Varili úžasne!!!

Ďalšou novinkou bola táborová mena „košáre“. Názov je odvodený od hlavnej budovy s~názvom „Košár“. Košáre sa dali získať prácou navyše, a potom si ich mohli vymeniť za niečo slané alebo sladké. Na naše prekvapenie sa najskôr minuli proteínové tyčinky \smiley. Aspoň vieme, čím tejto generácii urobíme najväčšiu radosť \smiley.

Pri záverečnom táboráku mali všetci možnosť povedať svoje hodnotenie tábora a na naše veľké potešenie sa pri otázke, čo sa im najviac páčilo, opakovali odpovede: téma, skupinky, chvály...
To je veľká Božia milosť, že v~ich srdciach rezonovala práve téma, ktorou bol list Kolosenským. A~našou túžbou bolo a stále je, aby sme my vedúci aj naši mladí videli, že v~Pánovi Ježišovi máme všetku plnosť a On stačí na všetko.
\autor{Miroslava Hovorková}


\clanok{Tábor v~Častej -- týždeň, ktorý stál za to}
Tábor v~Častej~2026 bol jednoducho vydarený. Zišlo sa nás okolo 80 a krásne
sa v~ňom stretli všetky generácie -- od najmenších až po tých, ktorí už majú za sebou nejeden tábor. A~práve táto pestrosť bola jednou z~vecí, ktoré robili celý týždeň výnimočným.

Prevažnú časť účastníkov tvorili pravidelní členovia BJB Palisády, no boli medzi nami aj takí, ktorí počas roka do cirkvi pravidelne nechodia.
Pre niektorých z~nich je tábor v~Častej dokonca jedinou príležitosťou počas roka stretnúť sa s~kresťanským spoločenstvom, stráviť spolu čas pri Božom slove, modlitbe, chválach a jednoducho byť medzi ľuďmi, ktorí svoju vieru žijú. Aj preto vnímam tábor ako oveľa viac než len príjemný letný týždeň.
Je to priestor, kde sa môžu budovať vzťahy, obnovovať spoločenstvo a možno aj nanovo objavovať cesta k~Bohu.

Atmosféra bola počas celého tábora veľmi pohodová. A~to napriek tomu, že náš program rozhodne nepatril medzi tie, pri ktorých si človek môže povedať: „Dnes nemám nič na pláne.“ Oficiálny deň začínal o~7.45~hod. rozcvičkou a končil o~22.00~hod. druhou večerou. Áno, druhou. Myslím, že po takomto programe sme sa na ňu všetci tešili.

Naše modlitby za dobré počasie, hladký priebeh tábora a hlavne za to, aby sme sa všetci vrátili domov bez zranení, boli vypočuté. Dokonca aj bazén nám bol tento rok k~dispozícii, čo sme počas veľkých horúčav ocenili možno viac ako inokedy.

Trochu nám bolo ľúto, že sa tento rok nikto neprihlásil na krst. Veríme však, že budúci rok túto krásnu príležitosť niekto využije.

Mojím osobným snom je, aby nás bolo každý rok viac a viac. Kapacita strediska nám však, žiaľ, pripomína, že aj dobré veci majú svoje limity.

Na príprave tábora sa podieľal celý tím ochotných ľudí. Každý priložil ruku
k~dielu podľa svojho obdarovania a vďaka tomu sa nikto nemusel „roztrhnúť“. Bola som veľmi vďačná, koľko ľudí úplne prirodzene prejavilo záujem zapojiť sa do príprav a prevziať zodpovednosť. Práve vďaka tomu mohol mať tábor takú pestrú podobu.

Každý deň sme mali duchovný program na tému Božia výzbroj, pričom sa
v~jednotlivých témach striedali bratia kazatelia. Ráno aj večer sme sa stretávali pri spoločných chválach, nechýbali programy pre deti z~besiedky, športové aktivity ako baseball, futbal, basketbal či pingpong, táborové hry boli tento rok naozaj vydarené, tvorivé dielničky boli milým spestrením a, samozrejme, nechýbali ani večerné táboráky s~pokojným zakončením dňa.

Táto podoba tábora sa rokmi postupne ustálila a ja by som chcela poďakovať všetkým generáciám pred nami, ktoré boli ochotné venovať svoj čas, energiu a ochotu, aby táto milá letná tradícia mohla pokračovať. Dnes z~nej čerpáme my a verím, že ju budeme odovzdávať ďalej.

Možno je to trochu ako v~biblickom obraze: „Púšťaj svoj chlieb po vode, lebo po mnohých dňoch ho nájdeš.“ (Kaz 11,1) To, čo niekto zasial pred rokmi svojím časom, službou a ochotou, sa dnes vracia v~podobe spoločných zážitkov, priateľstiev, služby a požehnania.

A~tak verím, že aj tábor v~Častej~2026 bol ďalším malým dielikom v~tejto krásnej reťazi. Niečo sme dostali od tých pred nami, niečo sme vytvorili spolu a niečo, dúfam, odovzdáme tým, ktorí prídu po nás.

Ďakujem všetkým, ktorí boli súčasťou tohto tábora -- či už svojou službou, modlitbami, ochotou, úsmevom, gitarou, loptou, alebo jednoducho tým, že prišli.

Teším sa na ďalší rok. A~možno nás bude opäť o~pár viac. Len budeme musieť presvedčiť aj Častú, aby trochu povolila svoju kapacitu. \smiley

\autor{za prípravný tím Diana Dzuriaková}


\clanok{Správy zo staršovstva}
Staršovstvo sa v~mesiaci august stretlo raz a na úvod sme sa zamysleli nad slovami apoštola Pavla „dávajte pozor na seba a na stádo“ ako pripomienkou, aby sme si všímali nielen svoje vlastné zápasy a radosti, ale aj zápasy a radosti ľudí v~zbore, a spoločne sa budovali v~Kristovi.

Venovali sme sa kandidatúre brata Jána Szőllősa na predsedu Rady BJB, ktorú staršovstvo odporúča, spolu s~návrhom Benjamína Uhrina za podpredsedu a Miroslava Tótha za člena KRK.

Prerokovali sme prípravu bohoslužieb od septembra~2026 -- novú kázňovú sériu na 1.~list Petrov, ktorá bude prebiehať od 13.~9. do začiatku adventu. Diskutovali sme aj o~zmysle „verša na mesiac“ a o~možnosti prepojiť spoločnú zborovú víziu so zborovým veršom na rok~2027. Padol návrh na obnovenie služby vítania počas bohoslužieb, ktorá by mala vytvárať priestor pre rozhovory a lepšie vzájomné spoznávanie sa.

Zaoberali sme sa aj prípravou inštalácie brata Dávida M.~Chuchúta za kazateľa, letnými aktivitami zboru ako aj zmluvou na prenájom budovy na Súľovskej a plánovaním spoločného vďakyvzdania so susednými zbormi BJB na 11.~10.~2026.
Tiež sme sa rozprávali o~sobáši Dávida a Juliany, o~ceste brata Jána Szőllősa na Ukrajinu na modlitebné raňajky do ukrajinského parlamentu a o~pripravovaných jesenných podujatiach -- Pochode za Ježiša a workshope Zmierenie Slovenska.

Prosíme vás o~modlitby za múdrosť pri rozhodovaní, za pokojné a pravdivé rozhovory, jednotu zboru a za všetky služby, ktoré máme ako zbor pred sebou.
\autor{Filip Barkóczi}


\clanok{Zbierky v~júli}
\table{lr}{
Na misiu           & 284~€\cr
Na investičný fond & 280~€\cr
}


\clanok{Zbierky v~auguste}
\table{lr}{
Na misiu           & 370~€\cr
Na investičný fond & 242~€\cr
}
\vskip1em

Pripomíname, že zbierky sa u~nás konajú pravidelne takto:
\vskip-1ex\begitems
* každú 2.~nedeľu v~mesiaci je zbierka venovaná misii a
* každú 4.~nedeľu je zbierka na investície.
\enditems

Prispieť je možné formou zbierky počas 4.~nedieľ v~mesiaci alebo bankovým prevodom na zborový účet, prípadne aj osobne s.~Ľubke Kohútovej.

Bankové spojenie: SK36 0900 0000 0000 1147 1836, SWIFT: GIBASKBX


\clanok{{\it ČO SA DEJE V~SEPTEMBRI}}


\clanok{Verš na mesiac}
V~septembri sa učíme nový verš, ktorý dostalo pre rok~2026 naše staršovstvo:
„Anjel Hospodinov táborí vôkol tých, čo sa ho boja a zachraňuje ich.“ (Ž 24,8)

V~nedeľu 6.9. sa uskutočnila slávnostná inštalácia za kazateľa brata Dávida M.~Chuchúta, po nej nasledoval spoločný obed v~Podunajských Biskupiciach.


\clanok{Očami Slávky}
Poznáte ten úsmev, ktorý ostáva na tvári ešte dlho po tom, ako ste zažili niečo pekné? Že sa pristihnete večer, či nasledujúci deň ako sa usmievate a myšlienkami ste v~predchádzajúcom dni...

Nedeľa bola pre mňa takým dňom. Dňom Božej blízkosti, požehnania aj ľudskej spolupatričnosti. Veľká rodina. Bola to sťaby svadba. Na jednej strane zbor, na druhej kazateľský pár. A~obaja prejavujú túžbu a vieru, že chcú spolu ťahať za jeden povraz aj spolu s~Najvyšším. Ani jeden si nechce robiť falošnú ilúziu o~ľudskej dokonalosti, ale spolu sa spoliehať, že pod Hospodinovou ochranou budeme kráčať spolu.

K~takýmto chvíľam patrí zvesť slova Božieho, spoločné chvály, čaro spevokolu, modlitby žehnania aj modlitby vďačnosti, sladké koláčiky a čas strávený spolu
v~spoločenstve.

Vďaka, Pane, že nám dávaš služobníkov, ktorí šíria Tvoju zvesť slovom aj piesňou, činmi aj budovaním vzťahov, počúvaním, viditeľnou aj neviditeľnou službou.

Taká bola pre mňa nedeľa. Vďaka, Nebeský Otče.
\autor{Slávka Kráľová}


\clanok{Uvítacia služba}
Od septembra sme obnovili uvítaciu službu pred bohoslužbami, ktorá vytvára priestor pre rozhovory a lepšie vzájomné spoznávanie sa. Chceme tak prispieť
k~tomu, aby sa u~nás každý cítil prijatý a ako doma. Veríme, že takáto služba pomôže nielen novým tváram medzi nami, ale posilní aj vzťahy v~celom našom spoločenstve.
Ak by ste mali záujem zapojiť sa, ozvite sa Filipovi Barkóczimu.


\clanok{Dotazník o~bohoslužbe}
V~minulosti náš zbor realizoval testy zamerané na duchovný stav zboru.
Z~výsledkov vyplýva, že jedna z~najviac problematických častí sú bohoslužby, ktoré boli málo duchovne inšpiratívne.
Od nedele 13.~9.~2026 začíname s~drobnými zmenami v~bohoslužbách.
Zároveň by sme vás chceli poprosiť o~vyplnenie krátkeho dotazníka, ktorý nám pomôže lepšie poznať aj Váš názor a skúsenosť.
Prosím, venujte dotazníku pár minút a pomôžte nám spoločne rásť aj v~tejto oblasti služby, ktorá neovplyvňuje len členov a priateľov nášho zboru, ale môže osloviť aj ľudí okolo nás.
\autor{staršovstvo BJB Palisády}
\vskip1em

Dotazník bol zaslaný v~e-mailových oznamoch.
Tí z~Vás, ktorí nemáte možnosť online prístupu a chcete vyplniť dotazník, prosím kontaktujte s.~K.~Kerekréty.


\clanok{Šarkaniáda s~Connectom}
V~sobotu 26.~9.~2026 sa môžete pripojiť ku Connect-ťákom pri púštaní šarkanov.
Stretnutie je od 11.00~hod. v~Novej dedinke na verejnej ploche futbalového ihriska (na konci cyklochodníka).


\n  2.	9.  Radislav NEMEC;
\n  5.	9.	Dušan	UHRIN;
\n 14.	9.	Štefan	SYNOVEC;
\n 16.	9.	Daniel	PLETT;
\n 19.	9.	Richard	HALAMÍČEK;
\n 21.	9.	Kvetoslava	MAĎAROVÁ;
\n 21.	9.	Miroslava	SIMONOVÁ;
\n 22.	9.	Viera	KOLÁROVSKÁ;
\n 25.	9.	Stanislav	KRÁĽ;
\narodeniny


\program{
\p  1 ; ut ;.;;.;;
\p  2 ; st ;.;;.;;
\p  3 ; št ;.;;.;;
\p  4 ; pi ; 17.30 ; Dorast ;.;;
\p  5 ; so ; 18.00 ; Mládež ;.;;
\p  6 ; ne ;  9.30 ; Bohoslužby (B.~Uhrin, inštalácia kazateľa);.;;
\p  7 ; po ; 18.00 ; Modlitebná skupinka ;.;;
\p  8 ; ut ; 15.00 ; Biblická hodina pre seniorov (P.~Pivka) ;.;;
\p  9 ; st ; 17.30 ; Stretnutie sestier ;.;;
\p 10 ; št ; 18.00 ; Biblická hodina (D.~M.~Chuchút) ;.;;
\p 11 ; pi ; 17.30 ; Dorast ;.;;
\p 12 ; so ; 18.00 ; Mládež ;.;;
\p 13 ; ne ;  9.30 ; Bohoslužby (D.~M.~Chuchút) ;.;;
\p 14 ; po ; 18.00 ; Modlitebná skupinka / skupinka „Rastieme v~Kristovi“ ;.;;
\p 15 ; ut ; 15.00 ; Biblická hodina pre seniorov (P.~Pivka) ;.;;
\p 16 ; st ;.;;.;;
\p 17 ; št ; 18.00 ; Biblická hodina (D.~M.~Chuchút) ;.;;
\p 18 ; pi ; 17.30 ; Dorast ;.;;
\p 19 ; so ; 18.00 ; Mládež ;.;;
\p 20 ; ne ;  9.30 ; Bohoslužby (J.~Szőllős) ;.;;
\p 21 ; po ; 18.00 ; Modlitebná skupinka ;.;;
\p 22 ; ut ; 15.00 ; Biblická hodina pre seniorov (P.~Pivka) ;.;;
\p 23 ; st ;.;;.;;
\p 24 ; št ; 18.00 ; Biblická hodina (D.~M.~Chuchút) ;.;;
\p 25 ; pi ; 17.30 ; Dorast ;.;;
\p 26 ; so ; 18.00 ; Mládež ;.;;
\p 27 ; ne ;  9.30 ; Bohoslužby (D.~M.~Chuchút) ;.;;
\p 28 ; po ; 18.00 ; Modlitebná skupinka / skupinka „Rastieme v~Kristovi“ ;.;;
\p 29 ; ut ; 15.00 ; Biblická hodina pre seniorov (P.~Pivka) ;.;;
\p 30 ; st ;.;;.;;
}

\tiraz
\bye
